%% Marcadores de revisão do orientador.
%% -------------------------------------------------------------------
%% Uso previsto: o orientador marca, em cor, o que revisou no texto de
%% Fabricio. A marcação e visual, nao altera o conteudo: o texto continua
%% identico quando a cor e desligada.
%%
%%   \pita{texto}                 texto revisado, na cor padrao (verde)
%%   \pita[azul]{texto}           texto revisado, em azul
%%   \pita[vermelho]{texto}       texto revisado, em vermelho
%%   \pita[caixa]{texto}          texto revisado em caixa independente
%%   \pita[caixa,vermelho]{texto} caixa com borda e texto em vermelho
%%   \pitacom[cor]{texto}{motivo} texto revisado + motivo
%%   \pitanota[cor]{motivo}       nota de rodape curta
%%   \begin{pitabloco}[cor] ... \end{pitabloco}   paragrafos inteiros
%%
%% Cores disponiveis: verde (padrao), azul e vermelho. A opcao `caixa'
%% pode ser combinada com uma cor, separada por virgula.
%%
%% A cor e desligada automaticamente na versao de entrega (opcao
%% `entrega' em main.tex) e pode ser desligado a mao com \pitacorfalse
%% antes de %% Marcadores de revisão do orientador.
%% -------------------------------------------------------------------
%% Uso previsto: o orientador marca, em cor, o que revisou no texto de
%% Fabricio. A marcação e visual, nao altera o conteudo: o texto continua
%% identico quando a cor e desligada.
%%
%%   \pita{texto}                 texto revisado, na cor padrao (verde)
%%   \pita[azul]{texto}           texto revisado, em azul
%%   \pita[vermelho]{texto}       texto revisado, em vermelho
%%   \pita[caixa]{texto}          texto revisado em caixa independente
%%   \pita[caixa,vermelho]{texto} caixa com borda e texto em vermelho
%%   \pitacom[cor]{texto}{motivo} texto revisado + motivo
%%   \pitanota[cor]{motivo}       nota de rodape curta
%%   \begin{pitabloco}[cor] ... \end{pitabloco}   paragrafos inteiros
%%
%% Cores disponiveis: verde (padrao), azul e vermelho. A opcao `caixa'
%% pode ser combinada com uma cor, separada por virgula.
%%
%% A cor e desligada automaticamente na versao de entrega (opcao
%% `entrega' em main.tex) e pode ser desligado a mao com \pitacorfalse
%% antes de %% Marcadores de revisão do orientador.
%% -------------------------------------------------------------------
%% Uso previsto: o orientador marca, em cor, o que revisou no texto de
%% Fabricio. A marcação e visual, nao altera o conteudo: o texto continua
%% identico quando a cor e desligada.
%%
%%   \pita{texto}                 texto revisado, na cor padrao (verde)
%%   \pita[azul]{texto}           texto revisado, em azul
%%   \pita[vermelho]{texto}       texto revisado, em vermelho
%%   \pita[caixa]{texto}          texto revisado em caixa independente
%%   \pita[caixa,vermelho]{texto} caixa com borda e texto em vermelho
%%   \pitacom[cor]{texto}{motivo} texto revisado + motivo
%%   \pitanota[cor]{motivo}       nota de rodape curta
%%   \begin{pitabloco}[cor] ... \end{pitabloco}   paragrafos inteiros
%%
%% Cores disponiveis: verde (padrao), azul e vermelho. A opcao `caixa'
%% pode ser combinada com uma cor, separada por virgula.
%%
%% A cor e desligada automaticamente na versao de entrega (opcao
%% `entrega' em main.tex) e pode ser desligado a mao com \pitacorfalse
%% antes de %% Marcadores de revisão do orientador.
%% -------------------------------------------------------------------
%% Uso previsto: o orientador marca, em cor, o que revisou no texto de
%% Fabricio. A marcação e visual, nao altera o conteudo: o texto continua
%% identico quando a cor e desligada.
%%
%%   \pita{texto}                 texto revisado, na cor padrao (verde)
%%   \pita[azul]{texto}           texto revisado, em azul
%%   \pita[vermelho]{texto}       texto revisado, em vermelho
%%   \pita[caixa]{texto}          texto revisado em caixa independente
%%   \pita[caixa,vermelho]{texto} caixa com borda e texto em vermelho
%%   \pitacom[cor]{texto}{motivo} texto revisado + motivo
%%   \pitanota[cor]{motivo}       nota de rodape curta
%%   \begin{pitabloco}[cor] ... \end{pitabloco}   paragrafos inteiros
%%
%% Cores disponiveis: verde (padrao), azul e vermelho. A opcao `caixa'
%% pode ser combinada com uma cor, separada por virgula.
%%
%% A cor e desligada automaticamente na versao de entrega (opcao
%% `entrega' em main.tex) e pode ser desligado a mao com \pitacorfalse
%% antes de \input{config/marcadores} (ou \pitacortrue para forcar).
%%
%% Nao usar dentro de titulos de secao nem de legendas: o hyperref
%% processa esses argumentos novamente e a cor vaza para o marcador do
%% PDF. Fora disso, o marcador funciona em qualquer lugar do texto.
%% -------------------------------------------------------------------

\RequirePackage{xcolor}

\definecolor{pitaverde}{HTML}{1A7F37}
\definecolor{pitaazul}{HTML}{1F6FEB}
\definecolor{pitavermelho}{HTML}{C00000}

%% Nomes curtos aceitos como parametro de cor.
\colorlet{verde}{pitaverde}
\colorlet{azul}{pitaazul}
\colorlet{vermelho}{pitavermelho}

%% \ifpitacor: verdadeiro enquanto as marcas de revisão estão ativas.
\newif\ifpitacor
\pitacortrue

%% Na versão de entrega as marcas somem sozinhas, sem depender de
%% lembrete: o texto revisado permanece, apenas sem a cor.
%% Todas as definicoes ficam dentro de \makeatletter porque os corpos
%% de \pita, \pitacom, \pitanota e pitabloco referenciam comandos
%% internos com `@'; fora deste bloco, o `@' teria catcode errado.
\makeatletter
\ifidp@entrega \pitacorfalse \fi

%% Interpreta a lista de opcoes do marcador: uma cor (verde, azul ou
%% vermelho) e/ou a palavra `caixa'. A cor padrao e verde.
\def\pita@caixanome{caixa}
\newif\ifpita@caixa
\newcommand*{\pita@cor}{verde}
\newcommand*{\pita@opcoes}[1]{%
  \pita@caixafalse
  \def\pita@cor{verde}%
  \@for\pita@item:=#1\do{%
    \ifx\pita@item\pita@caixanome
      \pita@caixatrue
    \else
      \let\pita@cor\pita@item
    \fi
  }%
}

%% Caixa independente do texto, sem recuo de paragrafo, com borda na
%% mesma cor do texto.
\newcommand*{\pita@caixaimprimir}[1]{%
  \par\noindent
  {\color{\pita@cor}%
    \fbox{\parbox{\dimexpr\linewidth-2\fboxsep-2\fboxrule\relax}{#1}}}%
  \par
}

%% Texto revisado. Opcoes: cor (verde padrao, azul, vermelho) e/ou `caixa'.
\newcommand*{\pita}[2][verde]{%
  \pita@opcoes{#1}%
  \ifpitacor
    \ifpita@caixa
      \pita@caixaimprimir{#2}%
    \else
      \textcolor{\pita@cor}{#2}%
    \fi
  \else
    #2%
  \fi}

%% Texto revisado com o motivo da intervenção logo em seguida.
\newcommand*{\pitacom}[3][verde]{\pita[#1]{#2}\ \pita[#1]{[#3]}}

%% Motivo em nota de rodapé.
\newcommand*{\pitanota}[2][verde]{\footnote{\pita[#1]{#2}}}

%% Trecho longo, com mais de um parágrafo, inteiramente na cor escolhida.
\newenvironment{pitabloco}[1][verde]%
  {\ifpitacor \color{#1}\fi \ignorespaces}%
  {\par}
\makeatother (ou \pitacortrue para forcar).
%%
%% Nao usar dentro de titulos de secao nem de legendas: o hyperref
%% processa esses argumentos novamente e a cor vaza para o marcador do
%% PDF. Fora disso, o marcador funciona em qualquer lugar do texto.
%% -------------------------------------------------------------------

\RequirePackage{xcolor}

\definecolor{pitaverde}{HTML}{1A7F37}
\definecolor{pitaazul}{HTML}{1F6FEB}
\definecolor{pitavermelho}{HTML}{C00000}

%% Nomes curtos aceitos como parametro de cor.
\colorlet{verde}{pitaverde}
\colorlet{azul}{pitaazul}
\colorlet{vermelho}{pitavermelho}

%% \ifpitacor: verdadeiro enquanto as marcas de revisão estão ativas.
\newif\ifpitacor
\pitacortrue

%% Na versão de entrega as marcas somem sozinhas, sem depender de
%% lembrete: o texto revisado permanece, apenas sem a cor.
%% Todas as definicoes ficam dentro de \makeatletter porque os corpos
%% de \pita, \pitacom, \pitanota e pitabloco referenciam comandos
%% internos com `@'; fora deste bloco, o `@' teria catcode errado.
\makeatletter
\ifidp@entrega \pitacorfalse \fi

%% Interpreta a lista de opcoes do marcador: uma cor (verde, azul ou
%% vermelho) e/ou a palavra `caixa'. A cor padrao e verde.
\def\pita@caixanome{caixa}
\newif\ifpita@caixa
\newcommand*{\pita@cor}{verde}
\newcommand*{\pita@opcoes}[1]{%
  \pita@caixafalse
  \def\pita@cor{verde}%
  \@for\pita@item:=#1\do{%
    \ifx\pita@item\pita@caixanome
      \pita@caixatrue
    \else
      \let\pita@cor\pita@item
    \fi
  }%
}

%% Caixa independente do texto, sem recuo de paragrafo, com borda na
%% mesma cor do texto.
\newcommand*{\pita@caixaimprimir}[1]{%
  \par\noindent
  {\color{\pita@cor}%
    \fbox{\parbox{\dimexpr\linewidth-2\fboxsep-2\fboxrule\relax}{#1}}}%
  \par
}

%% Texto revisado. Opcoes: cor (verde padrao, azul, vermelho) e/ou `caixa'.
\newcommand*{\pita}[2][verde]{%
  \pita@opcoes{#1}%
  \ifpitacor
    \ifpita@caixa
      \pita@caixaimprimir{#2}%
    \else
      \textcolor{\pita@cor}{#2}%
    \fi
  \else
    #2%
  \fi}

%% Texto revisado com o motivo da intervenção logo em seguida.
\newcommand*{\pitacom}[3][verde]{\pita[#1]{#2}\ \pita[#1]{[#3]}}

%% Motivo em nota de rodapé.
\newcommand*{\pitanota}[2][verde]{\footnote{\pita[#1]{#2}}}

%% Trecho longo, com mais de um parágrafo, inteiramente na cor escolhida.
\newenvironment{pitabloco}[1][verde]%
  {\ifpitacor \color{#1}\fi \ignorespaces}%
  {\par}
\makeatother (ou \pitacortrue para forcar).
%%
%% Nao usar dentro de titulos de secao nem de legendas: o hyperref
%% processa esses argumentos novamente e a cor vaza para o marcador do
%% PDF. Fora disso, o marcador funciona em qualquer lugar do texto.
%% -------------------------------------------------------------------

\RequirePackage{xcolor}

\definecolor{pitaverde}{HTML}{1A7F37}
\definecolor{pitaazul}{HTML}{1F6FEB}
\definecolor{pitavermelho}{HTML}{C00000}

%% Nomes curtos aceitos como parametro de cor.
\colorlet{verde}{pitaverde}
\colorlet{azul}{pitaazul}
\colorlet{vermelho}{pitavermelho}

%% \ifpitacor: verdadeiro enquanto as marcas de revisão estão ativas.
\newif\ifpitacor
\pitacortrue

%% Na versão de entrega as marcas somem sozinhas, sem depender de
%% lembrete: o texto revisado permanece, apenas sem a cor.
%% Todas as definicoes ficam dentro de \makeatletter porque os corpos
%% de \pita, \pitacom, \pitanota e pitabloco referenciam comandos
%% internos com `@'; fora deste bloco, o `@' teria catcode errado.
\makeatletter
\ifidp@entrega \pitacorfalse \fi

%% Interpreta a lista de opcoes do marcador: uma cor (verde, azul ou
%% vermelho) e/ou a palavra `caixa'. A cor padrao e verde.
\def\pita@caixanome{caixa}
\newif\ifpita@caixa
\newcommand*{\pita@cor}{verde}
\newcommand*{\pita@opcoes}[1]{%
  \pita@caixafalse
  \def\pita@cor{verde}%
  \@for\pita@item:=#1\do{%
    \ifx\pita@item\pita@caixanome
      \pita@caixatrue
    \else
      \let\pita@cor\pita@item
    \fi
  }%
}

%% Caixa independente do texto, sem recuo de paragrafo, com borda na
%% mesma cor do texto.
\newcommand*{\pita@caixaimprimir}[1]{%
  \par\noindent
  {\color{\pita@cor}%
    \fbox{\parbox{\dimexpr\linewidth-2\fboxsep-2\fboxrule\relax}{#1}}}%
  \par
}

%% Texto revisado. Opcoes: cor (verde padrao, azul, vermelho) e/ou `caixa'.
\newcommand*{\pita}[2][verde]{%
  \pita@opcoes{#1}%
  \ifpitacor
    \ifpita@caixa
      \pita@caixaimprimir{#2}%
    \else
      \textcolor{\pita@cor}{#2}%
    \fi
  \else
    #2%
  \fi}

%% Texto revisado com o motivo da intervenção logo em seguida.
\newcommand*{\pitacom}[3][verde]{\pita[#1]{#2}\ \pita[#1]{[#3]}}

%% Motivo em nota de rodapé.
\newcommand*{\pitanota}[2][verde]{\footnote{\pita[#1]{#2}}}

%% Trecho longo, com mais de um parágrafo, inteiramente na cor escolhida.
\newenvironment{pitabloco}[1][verde]%
  {\ifpitacor \color{#1}\fi \ignorespaces}%
  {\par}
\makeatother (ou \pitacortrue para forcar).
%%
%% Nao usar dentro de titulos de secao nem de legendas: o hyperref
%% processa esses argumentos novamente e a cor vaza para o marcador do
%% PDF. Fora disso, o marcador funciona em qualquer lugar do texto.
%% -------------------------------------------------------------------

\RequirePackage{xcolor}

\definecolor{pitaverde}{HTML}{1A7F37}
\definecolor{pitaazul}{HTML}{1F6FEB}
\definecolor{pitavermelho}{HTML}{C00000}

%% Nomes curtos aceitos como parametro de cor.
\colorlet{verde}{pitaverde}
\colorlet{azul}{pitaazul}
\colorlet{vermelho}{pitavermelho}

%% \ifpitacor: verdadeiro enquanto as marcas de revisão estão ativas.
\newif\ifpitacor
\pitacortrue

%% Na versão de entrega as marcas somem sozinhas, sem depender de
%% lembrete: o texto revisado permanece, apenas sem a cor.
%% Todas as definicoes ficam dentro de \makeatletter porque os corpos
%% de \pita, \pitacom, \pitanota e pitabloco referenciam comandos
%% internos com `@'; fora deste bloco, o `@' teria catcode errado.
\makeatletter
\ifidp@entrega \pitacorfalse \fi

%% Interpreta a lista de opcoes do marcador: uma cor (verde, azul ou
%% vermelho) e/ou a palavra `caixa'. A cor padrao e verde.
\def\pita@caixanome{caixa}
\newif\ifpita@caixa
\newcommand*{\pita@cor}{verde}
\newcommand*{\pita@opcoes}[1]{%
  \pita@caixafalse
  \def\pita@cor{verde}%
  \@for\pita@item:=#1\do{%
    \ifx\pita@item\pita@caixanome
      \pita@caixatrue
    \else
      \let\pita@cor\pita@item
    \fi
  }%
}

%% Caixa independente do texto, sem recuo de paragrafo, com borda na
%% mesma cor do texto.
\newcommand*{\pita@caixaimprimir}[1]{%
  \par\noindent
  {\color{\pita@cor}%
    \fbox{\parbox{\dimexpr\linewidth-2\fboxsep-2\fboxrule\relax}{#1}}}%
  \par
}

%% Texto revisado. Opcoes: cor (verde padrao, azul, vermelho) e/ou `caixa'.
\newcommand*{\pita}[2][verde]{%
  \pita@opcoes{#1}%
  \ifpitacor
    \ifpita@caixa
      \pita@caixaimprimir{#2}%
    \else
      \textcolor{\pita@cor}{#2}%
    \fi
  \else
    #2%
  \fi}

%% Texto revisado com o motivo da intervenção logo em seguida.
\newcommand*{\pitacom}[3][verde]{\pita[#1]{#2}\ \pita[#1]{[#3]}}

%% Motivo em nota de rodapé.
\newcommand*{\pitanota}[2][verde]{\footnote{\pita[#1]{#2}}}

%% Trecho longo, com mais de um parágrafo, inteiramente na cor escolhida.
\newenvironment{pitabloco}[1][verde]%
  {\ifpitacor \color{#1}\fi \ignorespaces}%
  {\par}
\makeatother